\chapter{Implementacija}

\section{Postavljanje Docker okruženja}

Svi servisi sustava, osim skripte za automatsku reakciju, izvode se u Docker kontejnerima \cite{docker_docs}. Napisane su dvije \texttt{docker-compose} datoteke, jedna za svaki
od dvaju ispitana podsustava. \texttt{docker-compose.yml} opisuje podsustav s bazom InfluxDB i agentom Telegraf, dok \texttt{docker-compose-elk.yml} opisuje podustav s bazom
Elasticsearch i agentom Filebeat. Obje konfiguracije dijele isti Zeek senzor i istu inačicu Grafane. Budući da oba podsustava koriste kontejner istog imena (\texttt{zeek-ids}) i isti
ulaz (\texttt{3000} za Grafanu), istovremeno pokretanje nije moguće, odnosno u bilo kojem danom trenutku se može izvoditi samo jedan podsustav, što je ujedno i preduvjet za poštenu usporedbu
u kasnijem poglavlju evaluacije.

Osjetljivi parametri poput korisničkih imena, lozinki, pristupnih tokena te naziva organizacije i spremnika \engl{bucket} — izdvojeni su iz \texttt{docker-compose} datoteka u zasebnu
\texttt{.env} datoteku okruženja, čime se konfiguracija odvaja od tajni što je standardna praksa. Struktura datoteke \texttt{.env} prikazana je u ispisu \ref{lst:env}.

\begin{lstlisting}[caption={Struktura \texttt{.env} datoteke (lozinke i pristupni tokeni skriveni)}, label={lst:env}]
INFLUXDB_USER=admin
INFLUXDB_PASSWORD=<lozinka>
INFLUXDB_ORG=dipl
INFLUXDB_BUCKET=zeek-logs
INFLUXDB_TOKEN=<token>

GRAFANA_USER=admin
GRAFANA_PASSWORD=<lozinka>
\end{lstlisting}

Pri prvom pokretanju InfluxDB se inicijalizira automatski putem varijabli \texttt{DOCKER\_INFLUXDB\_INIT\_*} u \texttt{setup} načinu rada \ref{lst:influx-init}, što
omogućuje da se bez ručne intervencije stvaraju korisnik, organizacija, spremnik i administratorski token. Time se izbjegava ručno postavljanje baze nakon svakog ponovnog pokretanja, 
što je veoma bitno za ponovljivost i održivost okruženja.

\begin{lstlisting}[caption={Automatska inicijalizacija InfluxDB okruženja iz \texttt{docker-compose.yml}}, label={lst:influx-init}]
environment:
  - DOCKER_INFLUXDB_INIT_MODE=setup
  - DOCKER_INFLUXDB_INIT_USERNAME=${INFLUXDB_USER}
  - DOCKER_INFLUXDB_INIT_PASSWORD=${INFLUXDB_PASSWORD}
  - DOCKER_INFLUXDB_INIT_ORG=${INFLUXDB_ORG}
  - DOCKER_INFLUXDB_INIT_BUCKET=${INFLUXDB_BUCKET}
  - DOCKER_INFLUXDB_INIT_ADMIN_TOKEN=${INFLUXDB_TOKEN}
\end{lstlisting}

Isti se token prosljeđuje Grafani kroz okolišnu varijablu kako bi se izvor podataka konfigurirao automatski pri pokretanju \engl{provisioning}, bez ručnog unosa pristupnih podataka
u sučelju kao što je vidljivo u ispisu \ref{lst:grafana-init}

\begin{lstlisting}[caption={Automatska inicijalizacija Grafane okruženja iz \texttt{docker-compose.yml}}, label={lst:grafana-init}]
environment:
  - GF_SECURITY_ADMIN_USER=${GRAFANA_USER}
  - GF_SECURITY_ADMIN_PASSWORD=${GRAFANA_PASSWORD}
  - GF_USERS_ALLOW_SIGN_UP=false
  # Prosljeduje token u provisioning YAML putem env varijable
  - INFLUXDB_TOKEN=${INFLUXDB_TOKEN}
  - INFLUXDB_ORG=${INFLUXDB_ORG}
  - INFLUXDB_BUCKET=${INFLUXDB_BUCKET}
\end{lstlisting}

Cjelokupno pokretanje, zaustavljanje i nadzor okruženja centralizira skripta \texttt{dipl.sh}, koja povezuje postavljanje pristupne točke i pokretanje kontejnera u jednu cjelinu. Postavljanje se
provodi korak po korak. Ako pokretanje kontejnera ne uspije, prethodno postavljena mrežna konfiguracija automatski se poništava \ref{lst:dipl-start}, čime se izbjegava
ostavljanje sustava u polupostavljenom stanju u kojem je pristupna točka aktivna, a nadzor nije.

\begin{lstlisting}[language=bash, caption={Atomarno pokretanje u \texttt{dipl.sh} (naredba \texttt{start})}, label={lst:dipl-start}]
if [ "$1" == "start" ]; then
  echo -e "${YELLOW}[1/2] Postavljanje AP okoline...${NC}"
  bash "$NET_SCRIPT" set
  if [ $? -ne 0 ]; then
    echo -e "${RED}Greska pri postavljanju AP-a. Prekidam.${NC}"
    exit 1
  fi
  echo -e "${GREEN}AP okolina postavljena.${NC}"
  echo -e "${YELLOW}[2/2] Pokretanje Docker servisa...${NC}"
  docker compose -f "$SCRIPT_DIR/docker-compose.yml" up -d
  if [ $? -ne 0 ]; then
    echo -e "${RED}Greska pri pokretanju Docker servisa. Resetiram AP...${NC}"
    bash "$NET_SCRIPT" reset
    exit 1
  fi
fi
\end{lstlisting}

Naredbe \texttt{stop}, \texttt{restart} i \texttt{status} upravljaju gašenjem, ponovnim pokretanjem i prikazom stanja servisa na sljedeći način: \texttt{sudo ./dipl.sh <naredba>}.

\section{Konfiguracija mreže, NAT-a i pristupne točke}

Potrebna mrežna konfiguracija postavlja se skriptom \texttt{net\_setReset\_rules.sh}, koja prima argument \texttt{set} ili \texttt{reset}. Uređaj raspolaže dvama sučeljima - žičanim
\texttt{enp2s0} prema pružatelju internetskih usluga i bežičnim \texttt{wlan1} u ulozi pristupne točke s adresom \texttt{10.0.1.1}. Bežično sučelje u ulozi pristupne točke ostvareno je 
USB bežičnim mrežnim prilagodnikom TP-Link TL-WN722N inačica v1.

Za rad u načinu pristupne točke prilagodnik mora zadovoljiti dva uvjeta. Njegov \textit{chipset} mora podržavati \engl{master} (AP) način rada, a pripadni upravljački
program mora biti dostupan kroz \texttt{nl80211} sučelje kojim \texttt{hostapd} upravlja prilagodnikom. Poželjno je i da je upravljački program dio jezgre čime se izbjegava
oslanjanje na vanjske, neslužbene upravljačke programe.

Inačica v1 zadovoljava sve navedeno. Temelji se na \textit{chipsetu} Atheros AR9271 s upravljačkim programom \texttt{ath9k\_htc}, koji je dio jezgre i izvorno podržava način
pristupne točke \cite{wikidevi_wn722n}. Kasnije inačice (v2 i v3) koriste \textit{chipset} Realtek RTL8188EUS, za koji podrška za način pristupne točke nije pouzdano dostupna kroz
jezgru, već zahtijeva ručno postavljanje neslužbenih upravljačkih programa. Stoga je upravo inačica v1 prikladna za ovu primjenu bez dodatnih zahvata.

Sučelje \texttt{wlan1} najprije se izuzima iz nadzora servisa \texttt{NetworkManager}, nakon čega mu se ručno dodjeljuje adresa i omogućuje prosljeđivanje paketa kao što je vidljivo
\ref{lst:nm-forward}. Izuzimanje je nužno jer bi \texttt{NetworkManager} inače pokušao samostalno upravljati sučeljem i ulazio u sukob s \texttt{hostapd}-om koji isto sučelje
postavlja u način rada pristupne točke.

\begin{lstlisting}[language=bash, caption={Izuzimanje sučelja i omogućavanje usmjeravanja}, label={lst:nm-forward}]
nmcli device set wlan1 managed no
ip link set wlan1 down
ip link set wlan1 up
ip addr add 10.0.1.1/24 dev wlan1
sysctl -w net.ipv4.ip_forward=1
\end{lstlisting}

Konfiguracija vatrozida i prevođenja adresa smještena su u zasebnu \texttt{nftables} tablicu \texttt{diplomski}, što omogućuje uklanjanje postavki jednom naredbom
bez utjecaja na ostatak sustava. Tablica sadrži pravilo maskiranja \engl{NAT} odlaznog prometa podmreže \texttt{10.0.1.0/24} na sučelju \texttt{enp2s0}, prikazano u ranijem
ispisu \ref{lst:nat}.

Uz pravilo maskiranja, ista tablica unaprijed priprema infrastrukturu za automatsku reakciju pomoću skupa \texttt{denylist} s vremenskim ograničenjem trajanja članova te dva lanca
koja odbacuju promet izvorišnih adresa iz tog skupa na \texttt{forward} i \texttt{input} kukama \engl{hooks} ref{lst:denylist}. Lanci su postavljeni na prioritet
\texttt{-10}, niži (a time hijerarhijski viši) od zadanih \texttt{ufw} lanaca, pa odbacivanje stupa na snagu prije ostalih pravila. Struktura blokade
definira se ovdje, jednom, dok skripta za reakciju opisana u kasnijoj sekciji samo dodaje adrese u postojeći skup.

\begin{lstlisting}[language=bash, caption={\texttt{denylist} i lanci za odbacivanje prometa}, label={lst:denylist}]
nft add set ip diplomski denylist { type ipv4_addr ; flags timeout ; timeout 30s ; }

nft add chain ip diplomski drop_fwd { type filter hook forward priority -10 ; }
nft add rule  ip diplomski drop_fwd ip saddr @denylist drop
nft add chain ip diplomski drop_in  { type filter hook input   priority -10 ; }
nft add rule  ip diplomski drop_in  ip saddr @denylist drop
\end{lstlisting}

Originalna politika vatrozida \texttt{ufw} odbacuje sav promet, pa se rad pristupne točke iznimkom dopušta najmanjim potrebnim skupom pravila \ref{lst:ufw} odnosno DHCP promet na
ulazima \texttt{67} i \texttt{68} te usmjeravanje s \texttt{wlan1} na \texttt{enp2s0}. Posljedica te je politike da će se dolazni neželjeni internetski promet odbaciti neovisno o ostalim slojevima.

\begin{lstlisting}[language=bash, caption={Dopuštenja \texttt{ufw} za rad pristupne točke}, label={lst:ufw}]
ufw allow in on wlan1 to any port 67 proto udp
ufw allow in on wlan1 to any port 68 proto udp
ufw route allow in on wlan1 out on enp2s0
\end{lstlisting}

Bežična pristupna točka realizirana je alatom \texttt{hostapd} \cite{hostapd_docs} na sučelju \texttt{wlan1}, uz zaštitu WPA2-PSK sa CCMP protokolom \ref{lst:hostapd}. Ključna
je postavka \texttt{ap\_isolate=0} kojom se onemogućuje izolacija klijenata. Pristupne točke uobičajeno izoliraju klijente kako međusobno ne bi komunicirali, no tada njihov
međusobni promet ne bi prolazio kroz točku nadzora; isključivanjem izolacije Zeek dobiva vidljivost i nad istok-zapad prometom unutar lokalne mreže.

\begin{lstlisting}[caption={Bitne postavke \texttt{hostapd} konfiguracije}, label={lst:hostapd}]
interface=wlan1
ssid=Dipl_test_FORBIDDEN
hw_mode=g
channel=6
wpa=2
wpa_key_mgmt=WPA-PSK
rsn_pairwise=CCMP
wpa_passphrase=<lozinka>
ap_isolate=0
\end{lstlisting}

Mrežne parametre postavlja \texttt{dnsmasq} \cite{dnsmasq_docs} \ref{lst:dnsmasq} tako da klijentima na sučelju \texttt{wlan1} dodjeljuje adrese iz raspona
\texttt{10.0.1.10}--\texttt{10.0.1.100}, zadaje usmjernik \texttt{10.0.1.1} te adrese vanjskih DNS poslužitelja \texttt{8.8.8.8} i \texttt{1.1.1.1}. U ovoj konfiguraciji
\texttt{dnsmasq} djeluje isključivo kao DHCP poslužitelj, dok se prevođenje adresa povjerava izravno vanjskim poslužiteljima, bez lokalnog prosljeđivanja. Time se ne gubi
nadzor nad DNS prometom, unatoč tome da su odredišni DNS poslužitelji izvan lokalne mreže, upiti svejedno prolaze kroz usmjernik i sučelje \texttt{wlan1}, gdje ih Zeek bilježi i pridodaje
pojedinom klijentu odnosno uređaju. Ova vidljivost vrijedi za klasične DNS upite preko UDP/53; klijenti koji koriste šifrirani DNS (DoH ili DoT) zaobilaze je jer Zeek takve upite vidi tek 
kao TLS odnosno QUIC promet prema vanjskom poslužitelju. Korišteni adresni prostor \texttt{10.0.1.0/24} dio je raspona \texttt{10.0.0.0/8} rezerviranog za privatne mreže prema RFC-u 1918 \cite{rfc1918}, 
dok dodjela adrese s završetkom \texttt{.1} usmjerniku slijedi uobičajenu, premda nestandardiziranu, praksu pri postavljanju lokalnih mreža.

\begin{lstlisting}[caption={Postavke DHCP-a i DNS prosljeđivanja u \texttt{dnsmasq}}, label={lst:dnsmasq}]
interface=wlan1
dhcp-range=10.0.1.10,10.0.1.100,255.255.255.0,24h
dhcp-option=option:router,10.0.1.1
dhcp-option=option:dns-server,8.8.8.8,1.1.1.1
\end{lstlisting}

Naredbom \texttt{sudo ./dipl.sh stop} u glavnom direktoriju cijela se konfiguracija poništava. Gase se \texttt{hostapd} i \texttt{dnsmasq}, briše se \texttt{nftables} tablica \texttt{diplomski},
uklanjaju se \texttt{ufw} iznimke, vraća se \texttt{ip\_forward} na 0 te se sučelje \texttt{wlan1} vraća pod nadzor \texttt{NetworkManagera}.

\section{Konfiguracija Zeek senzora}

Senzor je definiran kao servis \texttt{zeek} u obje \texttt{docker-compose} datoteke \ref{lst:zeekcmd}. Za razliku od ostalih servisa koji koriste izoliranu
\texttt{bridge} mrežu, Zeek se izvodi u načinu mreže domaćina \engl{host network mode} jer mora izravno pristupiti fizičkom sučelju \texttt{wlan1} dok u izoliranoj mreži kontejner
ne bi vidio promet domaćinskog sučelja. Zbog toga se njegovom kontejneru dodjeljuju ovlasti \texttt{NET\_ADMIN} i \texttt{NET\_RAW}, koje omogućuju otvaranje sirovih spojnih
točaka \engl{raw sockets} i hvatanje paketa na razini sučelja.

\begin{lstlisting}[caption={Definicija senzora u \texttt{docker-compose.yml}}, label={lst:zeekcmd}]
zeek:
  image: zeek/zeek:8.0.6
  network_mode: "host"
  cap_add: [ NET_ADMIN, NET_RAW ]
  volumes:
    - ./zeek-logs:/usr/local/zeek/logs
    - ./zeek/local.zeek:/usr/local/zeek/share/zeek/site/local.zeek:ro
    - ./zeek/intel.dat:/usr/local/zeek/share/zeek/site/intel.dat:ro
  command: zeek -C -i wlan1 local
\end{lstlisting}

Servis se pokreće naredbom \texttt{zeek -C -i wlan1 local}. Argument \texttt{-C} isključuje provjeru kontrolnih zbrojeva \engl{checksum} paketa (nužno jer mrežne kartice često
računaju zbrojeve sklopovski, pa ih jezgra Zeeku predaje kao naizgled neispravne), \texttt{-i wlan1} određuje sučelje za nadzor, a argument \texttt{local} učitava skriptu
\texttt{site/local.zeek}. Volumeni razdvajaju pisanje od čitanja tako što je \texttt{zeek-logs} montiran je s pravom pisanja kako bi senzor u njega zapisivao, dok su \texttt{local.zeek}
i \texttt{intel.dat} montirani isključivo za čitanje što je vidljivo iz dodatka \texttt{:ro} što znači \textit{read-only}.

Skripta \texttt{local.zeek} učitava bazne skenere protokola, okvir za obavijesti i proširenje \texttt{json-logs} pomoću kojega se zapisi pišu u JSON formatu prikladnom za izravno
strojno čitanje agenata. Uz Zeekov ugrađeni SSH \engl{bruteforce} detektor, u skripti su definirani i vlastite detektori. Ispis \ref{lst:zeek-detect} prikazuje detekciju skeniranja
portova. 

\begin{lstlisting}[caption={Detekcija skeniranja portova u \texttt{local.zeek}}, label={lst:zeek-detect}]
global port_scan_tracker: table[addr, addr] of set[port] &create_expire=5min;

event new_connection(c: connection)
{
    local src  = c$id$orig_h;
    local dest = c$id$resp_h;
    local dest_p = c$id$resp_p;

    if ([src, dest] !in port_scan_tracker)
        port_scan_tracker[src, dest] = set();
    add port_scan_tracker[src, dest][dest_p];

    if (|port_scan_tracker[src, dest]| >= 50)
        NOTICE([$note=Custom_Port_Scan, $src=src,
                $msg=fmt("Moguci port scan: %s kontaktirao %d razlicitih portova na %s",
                         src, |port_scan_tracker[src, dest]|, dest),
                $identifier=cat(src, dest), $suppress_for=30sec]);
}
\end{lstlisting}

Najprije nastaje tablica za praćenje ulaza koja se indeksira s parom [izvorišna adresa, odredišna adresa] a njihova vrijednost je skup portova koji su korišteni u komunikaciji.
Ukoliko je unutar pet minuta kontaktirano 50 ili više portova sa istim parom adresa, podiže se obavijest u \texttt{notice.log} i taj tip obavijesti se utišava na 30 sekundi kako upozorenja
ne bi zagušavala njihov zapisnik. Portovi se brišu iz skupa nakon pet minuta kako bi se očistila memorija. Analognim se pristupom prati i učestalost SSH konekcija prema usmjerniku.

Pokazatelji kompromitacije \engl{Indicators of Compromise} učitavaju se iz datoteke \texttt{intel.dat} putem Zeekovog Intel okvira \cite{zeek_docs}. Adrese se navode u datoteci, a
ne u skripti, čime se popis prijetnji može mijenjati bez izmjene detekcijske logike. Svi se zapisi razvrstavaju prema vrsti aktivnosti (\texttt{conn.log}, \texttt{dns.log},
\texttt{http.log}, \texttt{notice.log} i drugi) i pohranjuju u dijeljeni direktorij \texttt{zeek-logs}, kojem pristupaju podsustavi za pohranu opisani u nastavku.

\section{Telegraf/InfluxDB i Filebeat/Elasticsearch}

Podsustavi za pohranu kao ulaz koriste isti skup Zeekovih zapisa iz dijeljenog direktorija \texttt{zeek-logs}, montiranog isključivo za čitanje. Uzimaju se dva zapisa:
\texttt{conn.log}, koji nosi najveći dio prometa, i \texttt{notice.log}, koji sadrži upozorenja. Time se u obje baze unosi istovjetan ulaz, što je preduvjet usporedbe iz poglavlja
evaluacije. Razlike u rezultatu ovise jedino o pozadinskoj bazi i njenom agentu, a ne ulaznim podatcima.

U podsustavu s InfluxDB-om unos podataka obavlja Telegraf \cite{telegraf_docs}. Ulazni dodatak \texttt{inputs.tail} prati \texttt{conn.log} i \texttt{notice.log}, raščlanjuje svaki
redak kao JSON (format \texttt{json\_v2}) i upisuje ga u InfluxDB putem izlaza \texttt{influxdb\_v2}, kojem se pristupni token, organizacija i spremnik predaju putem varijabli okoliša.
Praćenje datoteka koristi metodu \texttt{poll} jer se zapisi nalaze na povezanom volumenu, gdje obavještavanje jezgre o promjenama datoteke (\texttt{inotify}) nije pouzdano.

Pri raščlanjivanju razdvajaju se oznake \engl{tags} od polja \engl{fields}. Oznake predstavljaju izvorišna i odredišna adresa, protokol, servis i stanje veze. InfluxDB ih indeksira, pa služe za
filtriranje i grupiranje u Grafani. Vrijednosti visoke ili jedinstvene raznolikosti, poput portova, količine prenesenih okteta i identifikatora veze \texttt{uid}, pohranjuju se kao
polja kako se ne bi nepotrebno umnažao broj vremenskih serija. Posebno specifično je raščlanjivanje Zeekovih ključeva oblika \texttt{id.orig\_h} s obzirom da \texttt{json\_v2} točku
tumači kao razdjelnik putanje, a u zapisu je ona dio samog naziva ključa, točka se u putanji mora izuzeti \ref{lst:telegraf-schema}.

\begin{lstlisting}[caption={Izuzimanje točke u nazivu ključa te podjela na oznaku i polje (\texttt{telegraf.conf})}, label={lst:telegraf-schema}]
[[inputs.tail.json_v2.tag]]
  path   = "id\\.orig_h"   # tocka je dio naziva kljuca, ne razdjelnik
  rename = "src_ip"

[[inputs.tail.json_v2.field]]
  path   = "id\\.resp_p"
  rename = "dst_port"
  type   = "int"
\end{lstlisting}

U podsustavu s Elasticsearchom ulogu agenta preuzima Filebeat \cite{filebeat_docs}. Svaki se Zeekov zapis prati zasebnim ulazom tipa \texttt{filestream} koji retke raščlanjuje kao
NDJSON, a pripadnu vrstu zapisa označava poljem \texttt{zeek\_type}. Procesor za vremenske oznake pretvara Zeekov \texttt{ts} (Unix vrijeme s decimalama) u standardno polje
\texttt{@timestamp}. Zapisi se potom šalju u Elasticsearch i pohranjuju u indeks nazvan prema vrsti i datumu, primjerice \texttt{zeek-conn-2026.05.29} \ref{lst:filebeat-input}.

\begin{lstlisting}[caption={Ulaz i odredišni indeks u \texttt{filebeat.yml}}, label={lst:filebeat-input}]
filebeat.inputs:
  - type: filestream
    id: zeek-conn
    paths: [ /zeek-logs/conn.log ]
    parsers:
      - ndjson: { keys_under_root: true }
    fields: { zeek_type: "conn" }

output.elasticsearch:
  hosts: [ "http://elasticsearch:9200" ]
  index: "zeek-%{[fields.zeek_type]}-%{+yyyy.MM.dd}"
\end{lstlisting}

Ostali tipovi zapisa (\texttt{dns.log, http.log, ssl.log, dhcp.log}) u konfiguraciji su pripremljeni, ali namjerno isključeni, kako bi oba podsustava
unosila istovjetan skup zapisa. Upravljanje predloškom indeksa i politikom životnog ciklusa indeksa \engl{Index Lifecycle Management, ILM} isključeno je, pa Elasticsearch
mapiranje polja stvara dinamički. Posljedice takvog mapiranja na zauzeće prostora razmatrati će se u poglavlju evaluacije.

\section{Grafana nadzorne ploče i usluživanje}

Grafana je za podsustave spremanja zajedničko vizualizacijsko sučelje \cite{grafana_docs}, čime se uklanja utjecaj načina prikaza na usporedbu pozadinskih baza. Izvor podataka ne
konfigurira se ručno, nego se učitava automatski pri pokretanju \engl{provisioning} iz datoteke opisa izvora, čime se izbjegava ručni unos pristupnih podataka i osigurava
ponovljivost. Za podsustav s InfluxDB-om opisuje se izvor tipa \texttt{influxdb} s Flux jezikom upita, kojem se organizacija, spremnik i token predaju okolišnim varijablama
\ref{lst:grafana-ds}, a za podsustav s Elasticsearchom opisuje se istovjetan izvor tipa \texttt{elasticsearch} nad uzorkom indeksa \texttt{zeek-*}. Budući da se podsustavi pokreću
zasebno, svaka instanca Grafane učitava isključivo svoj izvor podataka.

\begin{lstlisting}[caption={Automatsko učitavanje InfluxDB izvora podataka}, label={lst:grafana-ds}]
datasources:
  - name: InfluxDB
    type: influxdb
    access: proxy
    url: http://influxdb:8086
    uid: P951FEA4DE68E13C5      # fiksni uid radi podudaranja s nadzornom plocom
    jsonData:
      version: Flux
      organization: ${INFLUXDB_ORG}
      defaultBucket: ${INFLUXDB_BUCKET}
    secureJsonData:
      token: ${INFLUXDB_TOKEN}
\end{lstlisting}

Nadzorna ploča organizirana je u dvije skupine prikaza pregled prometa i sigurnosni pregled. Prikazi se temelje na shemi opisanoj u prethodnom odjeljku, odnosno na zapisima
\texttt{conn.log} i \texttt{notice.log} te pripadnim oznakama i poljima unutar tih zapisa.

Pregled prometa čine prikaz broja veza po minuti (koji broji zapise veza u kliznim jednominutnim prozorima), raspodjela protokola (kružni prikaz dobiven grupiranjem po oznaci
\texttt{proto}) te ljestvice najčešćih izvorišnih i odredišnih adresa. Primjer Flux upita, za prikaz broja veza po minuti, dan je u ispisu \ref{lst:flux-panel}; oznake se prije
zbrajanja uklanjaju kako bi se pojedinačne vremenske serije sažele u jednu vrijednost.

\begin{lstlisting}[caption={Flux upit prikaza broja veza po minuti}, label={lst:flux-panel}]
from(bucket: "zeek-logs")
  |> range(start: v.timeRangeStart, stop: v.timeRangeStop)
  |> filter(fn: (r) => r._measurement == "zeek_conn")
  |> filter(fn: (r) => r._field == "uid")
  |> drop(columns: ["src_ip", "dst_ip", "proto", "service", "conn_state"])
  |> aggregateWindow(every: 1m, fn: count, createEmpty: false)
  |> yield(name: "connections_per_minute")
\end{lstlisting}

Posebno je vrijedan prikaz najčešćih izvorišnih adresa. Budući da Zeek promet hvata na sučelju \texttt{wlan1} prije prevođenja adresa, izvorišne adrese ostaju stvarne adrese
uređaja u lokalnoj mreži (primjerice \texttt{10.0.1.94}), pa ih je u prikazu moguće i imenovati. Time prikaz izravno potvrđuje vrijednost odluke o mjestu nadzora iz poglavlja
arhitekture. Kada bi se hvatanje provodilo nakon prevođenja adresa, svi bi se izvorišni uređaji stopili u jednu adresu usmjernika i atribucija po uređaju bila bi izgubljena.

Sigurnosni pregled čine tablica upozorenja (najnovija upozorenja iz \texttt{zeek\_notice} s tipom i izvorišnom adresom), vremenski prikaz učestalosti upozorenja te tablica
sumnjivih veza. Potonja izdvaja veze čije stanje (\texttt{conn\_state}) odstupa od urednog dovršetka (\texttt{SF}) i grupira ih po stanju i izvorišnoj adresi, čime se ističu obrasci
poput velikog broja neuspjelih ili odbijenih veza s jednog izvora.

Istovjetna nadzorna ploča izrađena je i za podsustav s Elasticsearchom, s istim prikazima, ali s upitima prilagođenima Elasticsearch jeziku upita. Budući da se prikazi i njihov
raspored podudaraju, usporedba dvaju podsustava ne ovisi o načinu vizualizacije.

\section{Automatska reakcija: auto\_reaction.py i nftables}

Za razliku od ostalih komponenata, skripta za automatsku reakciju \texttt{auto\_reaction.py} izvodi se izravno na domaćinu, a ne u kontejneru, jer mijenja vatrozidna pravila domaćina
i za to zahtijeva administratorske ovlasti. Skripta povezuje detekciju, odnosno Zeekova upozorenja \cite{zeek_docs}, s odgovorom, tj. blokiranjem putem \texttt{nftables} kroz denylist.

Skripta prati najnovije retke datoteke \texttt{notice.log} naredbom \texttt{tail -F -n 0} gdje zastavica \texttt{-n 0} osigurava da se pri pokretanju zanemare postojeći zapisi, pa se
reagira isključivo na nova upozorenja, a ne na povijest zapisanu prije pokretanja.

Za svako novo upozorenje skripta čita njegov tip (\texttt{note}) i izvorišnu adresu. Izvorišna se adresa preuzima iz polja \texttt{src}, a ako ono ne postoji, iz polja
\texttt{id.orig\_h} \ref{lst:reaction-core}. Razlog je tomu što vlastiti detektori postavljaju \texttt{\$src} bez pridružene veze, dok upozorenja vezana uz konkretnu vezu
nose \texttt{id.orig\_h}.

\begin{lstlisting}[language=python, caption={Srž odlučivanja u \texttt{auto\_reaction.py}}, label={lst:reaction-core}]
note   = entry.get("note", "")
src_ip = entry.get("src") or entry.get("id.orig_h", "")

if note in TARGET_NOTICES and src_ip and src_ip not in WHITELIST:
    if src_ip not in BLOCKED_IPS:
        block(src_ip)   # nft add element ip diplomski denylist { src_ip }
\end{lstlisting}

Odgovor se pokreće isključivo za unaprijed određen skup upozorenja (\texttt{TARGET\_NOTICES}) kao pokušaji skeniranja portova, pogađanje SSH lozinke te podudaranja s pokazateljima
kompromitacije. Na taj način operator može smanjiti (ili povećati) broj upozorenja ovisno o situaciji i potrebama. Adresa usmjernika i adresa domaćina se nalaze u listi dozvoljenih 
(\texttt{WHITELIST = {"127.0.0.1", "10.0.1.1"}}) kako se ne bi mogla blokirati vlastita infrastruktura.

Blokada se provodi dodavanjem izvorišne adrese u skup \texttt{denylist} unutar tablice \texttt{diplomski}. Sam skup i lanci koji odbacuju promet definirani su ranije,
pri postavljanju mreže \ref{lst:denylist}. Skripta mijenja isključivo članstvo skupa, ne lance ni pravila. Zbog toga bilo kakav nenadani pad skripte ne ostavlja vatrozidna
pravila ne nadzirana već se jezgra brine o tome da se pobrišu nakon zadanog vremena.

Trajanje blokiranja ne određuje skripta, nego se očitava iz trenutne \engl{live} konfiguracije samog skupa (funkcija \texttt{read\_timeout}) kao jedinstvenog izvora istine. Istek pojedinog
elementa potom obavlja jezgra operacijskog sustava na temelju vremenskog ograničenja skupa, neovisno o skripti. Posljedica je otpornost na nasilni prekid. Prekine li se skripta
signalom \texttt{SIGKILL}, ponestankom memorije ili gubitkom napajanja, elementi svejedno isteknu sami, a cjelokupno se stanje uklanja brisanjem tablice. Skripta uz to vodi
vlastitu evidenciju blokiranih adresa isključivo radi izbjegavanja ponovnog dodavanja iste adrese, dok zasebna dretva uklanja istekle zapise iz te evidencije. Blokada se bilježi
tek ako je naredba \texttt{nft} uspješno izvršena, čime se izbjegava halucinirana evidencija nepostojeće blokade.

\section{Izazovi pri implementaciji}

Izgradnja sustava nije tekla bez zastoja. U nastavku se opisuju najznačajniji problemi uočeni tijekom implementacije i način na koji su riješeni, redom kako su se javljali kroz
slojeve sustava.

Već pri postavljanju kontejnerâ utvrđeno je da ranije korištena slika \texttt{zeekurity/zeek} više nije dostupna, pa je zamijenjena službenom slikom \texttt{zeek/zeek:8.0.6}.

Na mrežnom sloju prvi je izazov bio sukob oko upravljanja sučeljem. \texttt{NetworkManager} preuzimao je sučelje \texttt{wlan1} i gasio pristupnu točku koju podiže \texttt{hostapd}.
Problem je riješen izuzimanjem sučelja iz njegova nadzora naredbom \texttt{nmcli device set wlan1 managed no} prije konfiguracije pristupne točke \cite{arch_softap}.

Nakon što su \texttt{hostapd} i \texttt{dnsmasq} bili aktivni, a \texttt{wlan1} ispravno konfiguriran, klijenti i dalje nisu dobivali adresu. Pregledom \texttt{nft list ruleset}
utvrđeno je da \texttt{ufw} presreće DHCP pakete (ulaz \texttt{67/UDP}) i odbacuje ih prije nego dospiju do \texttt{dnsmasq}-a. Problem je otklonjen izrijekom dopuštenim DHCP prometom i
usmjeravanjem na sučelju \texttt{wlan1}, pravilima koja su potom ugrađena u skriptu \texttt{net\_setReset\_rules.sh} radi ponovljivosti.

Pokretanje \texttt{dnsmasq}-a isprva nije uspijevalo zbog zauzetosti ulaza \texttt{53} (\texttt{Address already in use}), koji je na povratnim adresama držao sistemski razlučivač
\texttt{systemd-resolved}. Budući da \texttt{dnsmasq} ionako treba slušati samo na \texttt{wlan1}, sukob je otklonjen direktivom \texttt{bind-interfaces} kojom \texttt{dnsmasq} sluša
isključivo na navedenom sučelju umjesto na svim adresama, bez gašenja sistemskog razlučivača. Uz to, prva inačica naziva mreže sadržavala je zagrade i razmake, koje \texttt{hostapd}
odbija kao nevaljane znakove u SSID-u, pa je naziv prilagođen.

Na sloju ingestije zapisi pri prvom pokretanju Telegrafa nisu dospijevali u InfluxDB. Uzrok je bilo pogrešno raščlanjivanje Zeekovih ključeva. Format \texttt{json\_v2} točku tumači
kao razdjelnik putanje, dok je u nazivima poput \texttt{id.orig\_h} ona dio samog ključa. Tek nakon izuzimanja točke u konfiguraciji (\texttt{id\textbackslash.orig\_h}) podatci su
ispravno raščlanjeni i upisani. Dodatno, datoteka \texttt{conn.log} ne nastaje pri pokretanju Zeeka nego tek pri prvoj zabilježenoj vezi, pa je Telegraf u međuvremenu prijavljivao
njezino nepostojanje; problem se rješava sam čim kroz mrežu prođe promet.

U pogledu detekcije isprva je za skeniranje portova učitana ugrađena politika \texttt{policy/misc/scan}. U korištenoj je inačici Zeeka ta politika zastarjela i prestala je
raditi, a prikladna zamjena uključena u Zeek instalaciju nije pronađena. Stoga je detekcija skeniranja preseljena u vlastiti detektor opisan ranije, čime je uklonjena ovisnost o
zastarjeloj i nepodržanoj politici.

Na sloju vizualizacije uočeno je da nadzorne ploče na čistom postavljanju ne pronalaze svoj izvor podataka. Uzrok je u tome što Grafana izvoru bez izrijekom zadanog identifikatora
dodjeljuje nasumičan \texttt{uid} \cite{grafana_docs}, koji se onda ne podudara s identifikatorom zapisanim u izvezenoj nadzornoj ploči. Problem je riješen ručnim postavljanjem polja
\texttt{uid} u opisu izvora podataka unutar \texttt{grafana/provisioning/datasources/influxdb.yaml} (i sličnoj putanji za Elasticsearch), čime se podudaranje 
osigurava pri svakom postavljanju.

Razlika između dvaju pozadinskih sustava očitovala se i pri vremenskim prikazima. Elasticsearch ograničava broj segmenata histograma (\texttt{search.max\_buckets}, zadano
\texttt{65536}); fiksni interval poput jedne minute nad širokim vremenskim rasponom premašuje tu granicu, što se kod InfluxDB-a ne događa. Prikazi su prilagođeni korištenjem
automatskog intervala koji se podešava prema odabranom rasponu. Ovo je ujedno konkretan primjer razlike u modelu dvaju sustava, podrobnije razmotrene u poglavlju evaluacije.

Naposljetku, pri opterećenju prometom generiranim alatom \texttt{t50} uočeno je da Zeek počinje gubiti pakete, što bilježi u zapisu \texttt{capture\_loss.log}, kada dotok premaši
približno pedeset tisuća paketa u sekundi. Za promet tipičan na lokalnoj mreži ciljanih razmjera to ne predstavlja ograničenje, no zabilježeno je kao gornja granica pouzdanog
hvatanja na korištenom sklopovlju.